\begin{abstract}
The mathematical attention mechanism has fundamentally transformed generative AI by enabling models to process vast contexts through dynamic information filtering. However, the theoretical link between this algorithmic structure and the behavioral "attention economy" of the modern Internet remains largely unexplored. In this work, we hypothesize that self-attention serves as a universal mechanism for capturing and distributing focus across both artificial neural networks and human-centric Internet products. We evaluate this hypothesis through a dual-pronged empirical analysis. First, we demonstrate that generative AI models utilize attention as a strict informational bottleneck, compressing uniform context distributions (Entropy 4.06 nats) into sparse, high-signal weights (Entropy 1.89 nats). Second, we model sequential Internet user behavior using a self-attention based recommender system, achieving a 68.5\% Hit Rate @ 5, which outperforms popularity-based content distribution by over 5.7$\times$. Our findings provide empirical evidence that the success of modern AI and Internet platforms is built upon the same mathematical foundation: an optimized mechanization of attention.
\end{abstract}
